\chapter{Maximum Power Point Tracking Algorithms}
\label{ch:mppt}

\section{All three algorithms are gradient estimators}

Section~\ref{sec:vanishing} established that the plant gain from control input
to power vanishes at the optimum, so no fixed linear controller can regulate to
it. What is needed instead is an estimate of the \emph{sign}, and ideally the
magnitude, of $\mathrm{d}P/\mathrm{d}V$. The three algorithms implemented here
differ only in how they obtain that estimate.

\begin{table}[htbp]
\centering
\caption{The three trackers viewed as gradient estimators.}
\label{tab:estimators}
\begin{tabular}{@{}p{0.22\textwidth}p{0.40\textwidth}p{0.28\textwidth}@{}}
\toprule
\textbf{Algorithm} & \textbf{How it estimates the gradient} & \textbf{Cost} \\
\midrule
Perturb and observe & Sign of $\Delta P$ between successive deliberate perturbations & Permanent limit cycle of $2\times$ the step \\
Incremental conductance & Direct evaluation of $\mathrm{d}I/\mathrm{d}V + I/V$ from successive samples & Same limit cycle, plus sensitivity to noise in a ratio \\
Extremum seeking & Demodulation of the response to a sinusoidal dither & Permanent dither of amplitude $A$ \\
\bottomrule
\end{tabular}
\end{table}

Every one of them pays a permanent steady-state cost, because every one of them
must keep exciting the plant to keep learning the gradient. That trade is
unavoidable and is quantified in Section~\ref{sec:tradeoff}.

\section{Perturb and observe}

\subsection{Principle}

The \term{perturb and observe} (P\&O) algorithm holds a direction $s \in \{+1,
-1\}$ and a fixed step $\Delta V$. Every sample period it moves the voltage
reference by $s\,\Delta V$, measures the resulting power, and reverses direction
if the power fell:

\begin{equation}
s_{k} = \begin{cases} -s_{k-1}, & P_k < P_{k-1} \\ \phantom{-}s_{k-1}, & \text{otherwise}\end{cases}
\qquad
V^\ast_{k+1} = V^\ast_k + s_k \Delta V .
\end{equation}

\subsection{The limit cycle}

At convergence the algorithm cannot stop; it has no way to know it has arrived.
It oscillates between three reference values, giving a steady-state limit cycle
of peak-to-peak amplitude exactly $2\Delta V$. Simulation confirms this
precisely: measured peak-to-peak values were \SI{1.00}{\volt},
\SI{1.99}{\volt}, \SI{2.99}{\volt}, \SI{3.98}{\volt}, \SI{5.96}{\volt},
\SI{7.98}{\volt} and \SI{11.95}{\volt} for steps of \SI{0.5}{\volt} to
\SI{6}{\volt}. This exactness is a useful validation that the voltage loop is
tracking the reference faithfully.

\subsection{The sampling constraint}

P\&O attributes the observed power change entirely to its own perturbation. Two
conditions must hold for that attribution to be valid:

\begin{enumerate}[leftmargin=1.6em]
  \item The voltage loop must have settled, or the measurement captures a
        transient. Chapter~\ref{ch:control} measured \SI{11.9}{\milli\second}
        settling, hence the \SI{20}{\milli\second} sample period.
  \item The irradiance must not have changed appreciably during the sample
        period. When it does, the attribution fails, which is the drift problem
        of Section~\ref{sec:drift}.
\end{enumerate}

\section{Incremental conductance}

\subsection{Principle}

\term{Incremental conductance} (INC)~\cite{hussein} rewrites the optimality
condition. Since $\mathrm{d}P/\mathrm{d}V = I + V\,\mathrm{d}I/\mathrm{d}V$,
setting it to zero gives

\begin{equation}
\frac{\mathrm{d}I}{\mathrm{d}V} = -\frac{I}{V},
\end{equation}

that is, the \emph{incremental} conductance equals the negative of the
\emph{instantaneous} conductance. The algorithm approximates the derivative by
$\Delta I / \Delta V$ from consecutive samples and moves the reference according
to the sign of the residual

\begin{equation}
r = \frac{\Delta I}{\Delta V} + \frac{I}{V} = \frac{1}{V}\frac{\mathrm{d}P}{\mathrm{d}V} .
\end{equation}

Because $V > 0$ always, $\mathrm{sign}(r) = \mathrm{sign}(\mathrm{d}P/\mathrm{d}V)$,
so the update rule is to increase the reference when $r > 0$ and decrease it
when $r < 0$. A tolerance band around zero is used to avoid chattering on noise.

\subsection{The case $\Delta V = 0$}

If the voltage has not changed between samples, the ratio is undefined. The
standard treatment, implemented here, examines $\Delta I$ instead: a positive
$\Delta I$ at constant voltage means the irradiance has risen, so the maximum
power point has moved to slightly higher voltage and the reference should
increase.

\subsection{Relationship to P\&O}

INC is often presented as fundamentally superior to P\&O. In the results of this
project the two are essentially identical under constant and slowly varying
irradiance, both reaching \SI{99.95}{\percent} at \SI{25}{\watt\per\square\metre\per\second},
and both showing the same $2\Delta V$ limit cycle. This should not be
surprising: with the same fixed step and the same sample period they are two
ways of computing the sign of the same quantity. INC's theoretical advantage,
that it evaluates a condition rather than differencing power, only materialises
when the two differ, which is under fast transients and in the presence of
measurement noise.

\section{Extremum seeking control}

\subsection{Structure}

The \term{extremum seeking} (ESC) tracker adds a sinusoidal dither to the
reference and extracts the gradient by correlation:

\begin{align}
V^\ast(t) &= \hat V(t) + A\sin(\omega t), \\
\xi(t) &= P(t) - \overline{P}(t) \quad\text{(washout high-pass)}, \\
\hat\gamma(t) &= \mathrm{LPF}\big[\xi(t)\sin(\omega t)\big] \;\approx\; \tfrac{A}{2}\frac{\mathrm{d}P}{\mathrm{d}V}, \\
\dot{\hat V}(t) &= K\,\hat\gamma(t) .
\end{align}

The washout filter removes the large DC component of power so that the
demodulator sees only the dither-induced variation. The low-pass filter after
demodulation removes the second-harmonic term. The integrator then performs
gradient ascent on the estimate.

\subsection{Parameter selection}

The dither frequency must satisfy competing constraints. It must be low enough
that the voltage control loop can actually follow the dither, so well below the
outer loop bandwidth of \SI{155}{\hertz} to \SI{223}{\hertz}. It must be high
enough that the dither is separable from genuine irradiance variation.
\SI{25}{\hertz} satisfies both. The washout and low-pass corners were set at
\SI{5}{\hertz}, a factor of five below the dither, and the amplitude at
\SI{3}{\volt}.

The adaptation gain $K$ was chosen by sweep, shown in Table~\ref{tab:escgain}.
The behaviour is exactly what averaging theory predicts: performance improves
with gain up to a point, then the estimator loop itself becomes unstable and the
steady-state excursion grows abruptly.

\begin{table}[htbp]
\centering
\caption{Extremum seeking gain sweep at \SI{1000}{\watt\per\square\metre}, starting from \SI{120}{\volt}. The true $\Vmp$ is \SI{157.80}{\volt}.}
\label{tab:escgain}
\begin{tabular}{@{}rrrr@{}}
\toprule
Gain $K$ & Static $\eta$ (\si{\percent}) & Estimate $\hat V$ (\si{\volt}) & Excursion (\si{\volt}) \\
\midrule
5   & 98.51 & 157.75 & 5.46 \\
20  & \textbf{99.21} & 157.75 & 5.51 \\
50  & 99.07 & 157.76 & 5.64 \\
100 & 98.62 & 157.79 & 5.89 \\
200 & 97.34 & 156.55 & 27.02 \\
400 & 97.67 & 161.31 & 27.06 \\
\bottomrule
\end{tabular}
\end{table}

$K = 20$ was adopted. Note that the internal estimate $\hat V$ converges to
\SI{157.75}{\volt} against a true value of \SI{157.80}{\volt}, an error of
\SI{0.03}{\percent}. ESC is the most \emph{accurate} of the three trackers; its
efficiency penalty comes entirely from the deliberate dither, not from bias.

\section{The fundamental trade-off}
\label{sec:tradeoff}

Every hill-climbing tracker faces the same trade: a larger perturbation acquires
the peak faster but oscillates further from it in steady state. This was
quantified by sweeping the P\&O step from \SI{0.5}{\volt} to \SI{6}{\volt},
measuring acquisition time from a \SI{120}{\volt} start and steady-state
efficiency separately, so that the two effects do not contaminate each other.

\begin{table}[htbp]
\centering
\caption{The P\&O speed against accuracy frontier at \SI{1000}{\watt\per\square\metre}.}
\label{tab:tradeoff}
\begin{tabular}{@{}rrrr@{}}
\toprule
Step (\si{\volt}) & Steady $\eta$ (\si{\percent}) & Limit cycle (\si{\volt}) & Acquisition (\si{\second}) \\
\midrule
0.5 & 99.9947 & 1.00  & 1.481 \\
1.0 & 99.9838 & 1.99  & 0.721 \\
1.5 & 99.9645 & 2.99  & 0.462 \\
2.0 & 99.9386 & 3.98  & 0.341 \\
3.0 & 99.7995 & 5.96  & 0.221 \\
4.0 & 99.6920 & 7.98  & 0.144 \\
6.0 & 99.4125 & 11.95 & 0.082 \\
\bottomrule
\end{tabular}
\end{table}

The relationship is cleanly monotone in both directions: an eighteen-fold
reduction in acquisition time costs \SI{0.58}{\percent} of steady-state
efficiency. A step of \SI{1.5}{\volt}, roughly \SI{1}{\percent} of $\Vmp$, was
adopted as a reasonable compromise and used for all comparisons.

\begin{figure}[htbp]
\centering
\includegraphics[width=\textwidth]{s04_mppt_comparison.png}
\caption[MPPT algorithm comparison]{Algorithm comparison through a \SI{1000}{} to \SI{500}{\watt\per\square\metre} irradiance step. Top left, voltage tracking; top right, power recovery; bottom left, steady-state behaviour showing the limit cycles and the extremum seeking dither; bottom right, the P\&O speed against accuracy frontier with each point labelled by its step size.}
\label{fig:mpptcompare}
\end{figure}

\section{Comparison under a step change}

Table~\ref{tab:stepcompare} compares all four tracker configurations through an
instantaneous irradiance step, the classic test.

\begin{table}[htbp]
\centering
\caption{Performance through a \SI{1000}{} to \SI{500}{\watt\per\square\metre} step.}
\label{tab:stepcompare}
\begin{tabular}{@{}lrrrr@{}}
\toprule
Tracker & $\eta$ (\si{\percent}) & RMS error (\si{\volt}) & Ripple (\si{\volt}) & Worst deficit (\si{\watt}) \\
\midrule
P\&O, fixed step        & 99.944 & 1.56 & 3.00 & 182.7 \\
P\&O, variable step     & \textbf{99.956} & 1.39 & \textbf{0.50} & \textbf{154.6} \\
Incremental conductance & 99.944 & 1.56 & 3.00 & 182.7 \\
Extremum seeking        & 99.551 & 4.88 & 5.95 & 194.6 \\
\bottomrule
\end{tabular}
\end{table}

Under this test the variable-step P\&O looks best: it shrinks its step near the
peak, achieving a \SI{0.5}{\volt} limit cycle, and enlarges it during the
transient. Extremum seeking looks worst, penalised by its \SI{6}{\volt}
peak-to-peak dither.

Chapter~\ref{ch:benchmark} shows that this ranking inverts completely under
irradiance \emph{ramps}, which is why a step test alone is an inadequate
benchmark and why the EN 50530 ramp profiles exist.

\section{The drift problem}
\label{sec:drift}

An instantaneous step is, paradoxically, an easy test for P\&O. The irradiance
is constant both before and after, so the power change the algorithm observes is
genuinely caused by its own perturbation except in the single sample that
straddles the step.

A \emph{ramp} is much harder. Suppose irradiance is rising at
\SI{100}{\watt\per\square\metre\per\second} and the sample period is
\SI{20}{\milli\second}. Between samples the irradiance rises by
\SI{2}{\watt\per\square\metre}, which at full sun adds roughly \SI{5}{\watt} to
the available power. A \SI{1.5}{\volt} perturbation near the flat peak changes
power by far less than that. So $\Delta P > 0$ regardless of which way the
algorithm perturbed, it concludes its last move was correct, and it keeps going
in that direction. If that direction happens to be wrong, it marches steadily
away from the peak. This is the well-documented \term{drift} or
\term{confusion} phenomenon~\cite{femia}.

Whether it occurs depends on the phase of the perturbation when the ramp starts,
which is why it is intermittent and easy to miss in a single simulation. The
next chapter observes it directly, and shows that a naive variable-step rule,
which was the best performer under the step test, makes it dramatically worse.
